%==========================================================
% No enunciado, a introdução não é numerada: o brief é a Secção 1.
\section*{Introdução}\label{sec:introducao}
\addcontentsline{toc}{section}{Introdução}

\subsection*{O projeto \projeto}

O \projeto{} é um agente de viagens baseado em inteligência artificial,
pensado, numa primeira fase, para jovens com orçamentos limitados
mas posteriormente para qualquer utilizador que queira simplificar o processo de planeamento de viagens.
Ainda pretendemos avaliar o potencial do \projeto{} em atender às novas necessidades de agências de viagens já
inseridas no mercado.

Na Entrega~1 estudámos como esse trabalho é feito hoje, com duas entrevistas contextuais: um jovem que organizou uma
viagem de grupo e um agente de viagens de balcão. Os modelos de trabalho mostraram que o problema não é só encontrar
opções baratas, mas saber quanto a viagem custa de facto, coordenar o grupo e o dinheiro e confiar no que a IA sugere.
Por isso a nossa visão deixou de ser um agente que propõe o itinerário sozinho e passou a ser um quadro que o
utilizador monta e que a IA vigia. Nesta entrega, o \projeto{} também não faz reservas nem processa pagamentos: o
organizador escolhe a proposta combinada e reserva no site de cada fornecedor.

Esta entrega descreve o sistema a construir a partir desses resultados.

\subsection*{Contexto}

Este projeto foi desenvolvido na unidade curricular de \uc{}, no Mestrado em Engenharia Informática do DEI da Universidade de Coimbra.
Não há um cliente externo: a ideia do \projeto{} partiu da equipa. Por isso, o \textit{brief} (Secção~\ref{sec:brief}) é a
ideia inicial com que a equipa propôs o projeto.

\subsection*{Estrutura do documento}

A Secção~\ref{sec:brief} apresenta o \textit{brief}. A Secção~\ref{sec:storymap} contém o \textit{story map} e as histórias de
utilizador. A Secção~\ref{sec:casosdeuso} apresenta o diagrama de contexto, o diagrama de casos de uso, os três casos de uso
detalhados e uma reflexão sobre um caso de uso gerado por IA. A Secção~\ref{sec:bdd} contém os cenários BDD, a
Secção~\ref{sec:ddd} o modelo de domínio, a Secção~\ref{sec:reflexao} a reflexão sobre o uso de IA e a
Secção~\ref{sec:conclusao} as conclusões.
